\scanpage{50}
% archive sheet number: 49

$X_{\unsure{\mathrm{car}{\geq}0|\mathrm{car}}}\simeq\widehat{R_{\mathrm{car}}}$, où $R_{\mathrm{car}}$ est la catégorie somme amalgamée des $\unsure{K_p|R_p}$ (tp fini \deleted{\unsure{sur $\mathbb Z/p$|corps fini}}) sur $\varnothing\longrightarrow R_p$ pour tt $p$, i.e. des schémas \deleted{\unsure{de|sur}} qui sont de type fini sur un corps premier (y inclus $\mathbb Q$) \deleted{\unsure{ayant un|ayant des} \illegible{2}}

21) \uline{Anneaux à car. avec propriétés précisées}

\deleted{\unsure{21)|22)}} \uline{Re NB} On peut repêcher maintenant toutes les notions considérées \inserted{dans le cas de car. donnée en propr\textsuperscript{té} précise} \deleted{ci plus haut, et donner pour} tt $p$ (fini ou infini) un sous-topos \deleted{\ensuremath{\unsure{X|Y}}} $Z(p)$ de $\widehat{R_p}$, on trouvera comme topos modulaire correspondant le topos somme de ceux des $Z(p)$ tels que $\varnothing_p\notin Z(p)$ i.e. l'anneau nul \deleted{\unsure{n'en|n'est}} n'est pas \deleted{dans} du type $Z(p)$, et de la somme amalgamée sur \deleted{la catégorie} \inserted{le topos} ponctuelle des autres $Z(p)$ [munis de leur sous-topos ouvert ponctuel $\varnothing_p$]. Il y a assez de pts ssi chacun des $Z(p)$ \deleted{\unsure{en|on}} a assez

22) Anneau quasi \deleted{\unsure{irrédu|indé}} connexe de rang $\leq n$ (dont le spectre a au plus $n$ comp connexes \deleted{\unsure{$\}$|?}}, donc qui est somme d'au plus $n$ schémas connexes)

\deleted{A} \uline{Déf} On peut trouver loc. une décomposition en produit $A_1\times A_2\times\cdots\times A_n$, avec les $A_i$ connexes $\Longleftrightarrow$ on peut trouver loc. \deleted{d'un produit} des projecteurs $e_i$ ($1\leq i\leq n$) orthogonaux de somme $1$, avec $A_j=A/\sum_{j\ne i}e_j A$ connexe pour tt $1\leq j\leq n$.

\diagram{\begin{tikzcd}[ampersand replacement=\&,row sep=small]
 \Pi_n\subset A^n\arrow[d,hookrightarrow] \\
 e_1,\ldots,e_n\ \text{avec}
\end{tikzcd}}

\[
 \left\{\begin{aligned}
 e_i^2&=e_i\quad\forall i\\
 e_i e_j&=0\quad\forall i,j\quad i\ne j\\
 \sum e_i&=1
 \end{aligned}\right.
\]

\diagram{\begin{tikzcd}[ampersand replacement=\&,row sep=small]
 A\times\Pi_n\arrow[d]\& A_j=A_{\Pi_n}/\displaystyle\sum_{j\ne i}e_j A_{\Pi_n}\arrow[dl]\\
 \deleted{\ensuremath{\unsure{A|p}}}\Pi_n\arrow[d,leftrightarrow]\&\\
 {}\&
\end{tikzcd}}

\diagram{\begin{tikzcd}[ampersand replacement=\&,row sep=small]
 \Pi_n\amalg\Pi_n\arrow[r,"{i_{n,j}}"]\arrow[dr,dash]\&\Pi(A_j)\arrow[d,dash]\quad\text{(proj. de $A_j$)}\\
 \&\Pi_n
\end{tikzcd}}

\unsure{2|?}\qquad soit $\Pi_{nj}\subset R_n$ plus grand ss-objet pour lequel \unsure{$i_{nj}$ un épi.|$i_{nj}$ une épi} commut\textsuperscript{t} avec chgt de base \deleted{dans} \uline{$\mathcal E$}
